\documentclass{article}

\usepackage{booktabs}
\usepackage{tabularx}

\title{Development Plan\\\progname}

\author{\authname}

\date{}

\input{../Comments.text}
\input{../Common.text}

\begin{document}

\maketitle

\begin{table}[hp]
\caption{Revision History} \label{TblRevisionHistory}
\begin{tabularx}{\textwidth}{llX}
\toprule
\textbf{Date} & \textbf{Developer(s)} & \textbf{Change}\\
\midrule
2026/09/25 & Jake Finlay & Completed Section 10: Proof of Concept Demonstration Plan \\
\bottomrule
\end{tabularx}
\end{table}

\newpage{}

\wss{Put your introductory blurb here.  Often the blurb is a brief roadmap of
what is contained in the report.}

\wss{Additional information on the development plan can be found in the
\href{https://gitlab.cas.mcmaster.ca/courses/capstone/-/blob/main/Lectures/L02b_POCAndDevPlan/POCAndDevPlan.pdf?ref_type=heads}
{lecture slides}.}

\section{Confidential Information?}

\wss{State whether your project has confidential information from industry, or
not.  If there is confidential information, point to the agreement you have in
place.}

\wss{For most teams this section will just state that there is no confidential
information to protect.}
\section{IP to Protect}

\wss{State whether there is IP to protect.  If there is, point to the agreement.
All students who are working on a project that requires an IP agreement are also
required to sign the ``Intellectual Property Guide Acknowledgement.''}

\section{Copyright License}

\wss{What copyright license is your team adopting.  Point to the license in your
repo.}

\section{Team Meeting Plan}

\wss{This section is for meetings between the team members, and meetings between 
the team and their supervisor/advisor (as appropriate).}

\wss{How often will the team meet? where? when?}

\wss{If the meeting is a physical location (not virtual), out of an abundance of
caution for safety reasons you shouldn't put the location online}

\wss{To help schedule meetings, you can use the lecture/tutorial time when we do 
not have anything else scheduled.}

\wss{How often will you meet with your supervisor/advisor?  when?  where?}

\wss{Will meetings be virtual?  At least some meetings should likely be
in-person.}

\wss{How will the meetings be structured?  There should be a chair for all meetings.  
There should be an agenda for all meetings.}

\section{Team Communication Plan}

\wss{Issues on GitHub should be part of your communication plan.}

\section{Team Member Roles}

\wss{You should identify the types of roles you anticipate, like notetaker,
leader, meeting chair, reviewer.  Assigning specific people to those roles is
not necessary at this stage.  In a student team the role of the individuals will
likely change throughout the year.}

\section{Workflow Plan}

\begin{itemize}
	\item How will you be using git, including branches, pull request, etc.?
	\item How will you be managing issues, including template issues, issue
	classification, etc.?
  \item Use of CI/CD
\end{itemize}

\section{Project Decomposition and Scheduling}

\begin{itemize}
  \item How will you be using GitHub projects?
  \item Include a link to your GitHub project
\end{itemize}

\wss{How will the project be scheduled?  This is the big picture schedule, not
details. You will need to reproduce information that is in the course outline
for deadlines.}

\section{Proof of Concept Demonstration Plan}

The proof of concept (PoC) demonstration will focus on one complete ``vertical
slice'' of the application instead of covering many procedures at once. The team
will demonstrate the following on both a physical iOS device and a physical
Android device:

\begin{enumerate}
  \item A child-friendly walkthrough of one common procedure (receiving eye
  drops), shown through a guide character with narration, simple language, and
  illustrations or animations.
  \item An interactive vision testing practice activity, where the child matches
  picture-based symbols similar to those used on paediatric vision charts.
  \item Navigation between these activities that a child could complete with
  little or no help from a parent.
\end{enumerate}

\noindent This slice was chosen because it exercises every major risk to the project.
The accuracy of clinical content, whether young children can actually use and
engage with the application, and whether rich media can be delivered reliably
on both mobile platforms.

\subsection{Risk 1: Clinical Accuracy and Age-Appropriateness of Content}

The value of the application depends on children getting an accurate preview of
what will happen at their appointment. If a procedure is shown incorrectly (for
example, steps in the wrong order, leaving out that eye drops may sting, or
equipment that looks different from what is in the clinic), the child's
expectations will not match reality. This could increase anxiety, which is the
opposite of the project's goal. No team member has clinical training, so there
is a risk that the team cannot produce content that is both clinically accurate
and understandable for young children. There is also a risk that getting
content reviewed by a clinician takes too long to cover the full set of planned
procedures.

\subsubsection*{What the PoC will demonstrate:}
\begin{itemize}
  \item The eye drop walkthrough (script and visuals) will be reviewed by a
  paediatric ophthalmology clinician before the demonstration. The team will
  record the corrections requested and how many review rounds were needed.
  \item The narration will use short sentences and familiar words, and will be
  honest about sensations (e.g., drops may feel cold or sting briefly) while
  staying reassuring.
  \item The team will measure how long one procedure takes to go from draft to
  clinician-approved. This will be used to estimate how many procedures are
  realistic for the final product.
\end{itemize}

\subsubsection*{Success criteria:}
The clinician confirms the walkthrough accurately
reflects what a child experiences, with no outstanding corrections, and the
review turnaround is short enough to cover the remaining planned procedures
(e.g., vision testing and slit-lamp examination).

\subsection{Risk 2: Child Engagement and Usability}

The target users are young children, many of whom cannot read yet, have short
attention spans, and are still developing fine motor skills. If they cannot
navigate the application on their own or lose interest partway through, they
will not benefit from the preparation. Designing for children is different from
designing for adults and requires audio-first instructions, large touch targets,
and interactions that are forgiving of mistakes. A related risk is that testing
directly with children may require research ethics approval, which may not be
in place by the PoC date.

\subsubsection*{What the PoC will demonstrate:}
\begin{itemize}
  \item All instructions in the vertical slice are narrated, so no step
  depends on the user reading text.
  \item Touch targets are sized for small fingers, and incorrect taps in the
  vision practice activity lead to encouragement instead of failure.
  \item Informal walkthroughs will be run with people unfamiliar with the
  application (e.g., parents or peers outside the team). The team will record
  whether they complete the walkthrough and activity without verbal help, how
  long it takes, and where they get stuck. If ethics approval is in place, a
  supervised session with a child will also be run.
  \item The team will confirm what approval is required for testing with
  children and a realistic timeline for obtaining it.
\end{itemize}

\subsubsection*{Success criteria:}
Testers complete both activities without verbal
help, the full slice takes about five minutes or less so that it fits a young
child's attention span, and the team has a clear path to testing with children
before the final evaluation.

\subsection{Risk 3: Cross-Platform Delivery of Rich Media}

The application must run on iOS and Android, on both phones and tablets, while
playing animations, audio narration, and interactive touch content. Families may
use older or lower-end devices, and the application may be opened in hospital
waiting rooms where connectivity is poor. As more procedures are added, media
files could make the application large to download. With a small team, keeping
one codebase working well on both platforms is a technical risk.

\subsubsection*{What the PoC will demonstrate:}
\begin{itemize}
  \item The same build of the vertical slice running on at least one iOS device
  and one Android device, including at least one older or lower-spec device.
  \item Narration and animations playing smoothly and in sync on each device.
  \item The slice working with no network connection (airplane mode) after
  installation.
  \item The install size of the slice, used to estimate the total size of the
  application once all planned procedures are included.
\end{itemize}

\subsubsection*{Success criteria:}
The slice behaves the same on both platforms with no
noticeable lag or audio/visual desynchronization, works fully offline, and the
estimated full application size is reasonable for a family to download.

\subsection{Consequences of an Unsuccessful PoC}

If the PoC does not sufficiently address these risks, the consequences for the
project are as follows.

\subsubsection*{If clinical accuracy (Risk 1) is not resolved:}
\begin{itemize}
  \item The team cannot claim the content is trustworthy. Inaccurate content
  could mislead children and increase their anxiety, and clinicians would be
  unlikely to recommend the application to families.
  \item If clinician review is too slow, fewer procedures can be covered in the
  final product (e.g., only eye drops and vision testing, without the slit-lamp
  examination), leaving parts of the appointment unprepared for.
\end{itemize}

\subsubsection*{If child engagement and usability (Risk 2) are not resolved:}
\begin{itemize}
  \item If children cannot use the application on their own, it becomes a
  parent-led tool. This would require rethinking the interaction design and may
  make the application less engaging for the child.
  \item If approval to test with children cannot be obtained in time, the final
  evaluation would have to rely on adult testers and clinician judgement. This
  would weaken any claims that the application actually helps children feel
  more prepared.
\end{itemize}

\subsubsection*{If cross-platform delivery (Risk 3) is not resolved:}
\begin{itemize}
  \item Supporting only one platform would exclude families who use the other,
  reducing how accessible the application is.
  \item Poor performance on older devices, or a very large download, would
  create barriers for some families. The team may need to reduce the richness
  of the media (e.g., static illustrations instead of animations), which could
  make the application less engaging for children.
\end{itemize}

\noindent Whether or not the PoC succeeds, its results will be used to set the number of
procedures covered, the interaction design, and the media approach.

\section{Expected Technology}

\wss{What programming language or languages do you expect to use?  What external
libraries?  What frameworks?  What technologies.  Are there major components of
the implementation that you expect you will implement, despite the existence of
libraries that provide the required functionality.  For projects with machine
learning, will you use pre-trained models, or be training your own model?  }

\wss{The implementation decisions can, and likely will, change over the course
of the project.  The initial documentation should be written in an abstract way;
it should be agnostic of the implementation choices, unless the implementation
choices are project constraints.  However, recording our initial thoughts on
implementation helps understand the challenge level and feasibility of a
project.  It may also help with early identification of areas where project
members will need to augment their training.}

Topics to discuss include the following:

\begin{itemize}
\item Specific programming language
\item Specific libraries
\item Pre-trained models
\item Specific linter tool (if appropriate)
\item Specific unit testing framework
\item Investigation of code coverage measuring tools
\item Specific plans for Continuous Integration (CI), or an explanation that CI
  is not being done
\item Specific performance measuring tools (like Valgrind), if
  appropriate
\item Tools you will likely be using?
\end{itemize}

\wss{git, GitHub and GitHub projects should be part of your technology.}

\section{Use of AI and LLMs}

\wss{This section is a continuation of the previous section.  The role of LLMs
in your project as important enough to warrant their own section.  
How will your team be using LLMs for code and documentation development?  
What tools will you be using?  How will you verify the output of your LLMs?}

\section{Coding Standard}

\wss{What coding standard will you adopt?}

\newpage{}

\section{Appendix --- Generative AI Use Disclosure}

\wss{ChatGPT (version 5) was used to brainstorm ideas for the template for this disclose section.}

\wss{This section is about the use of AI to write this document, not your
planned use of AI for your project.  You discuss that topic in one of the
sections in the body of this report.}

\subsection{AI Tools Used}

\wss{Give the name of the AI tool(s) used and the version.}

\subsection{How and Where Used}

\wss{Explain what the tool(s) were used for.  Examples include: brainstorming,
explaining a technical concept, reviewing the document, generating or modifying
text, generating or modifying diagrams, identifying errors or omissions or
inconsistencies.}

\wss{For each of the uses, identify which parts of the document were affected.}

\wss{You do not need to report every interaction, but you should disclose the 
uses that materially contributed to your work}

\subsection{Verification of AI-Generated Content}

\wss{\begin{itemize}
    \item Describe how the team checked AI-generated or AI-assisted material
    before including it in the document.
    \item In particular, verify factual claims, technical assertions,
    requirements, calculations, references, and other information for which an
    AI system may produce incorrect or fabricated results.
    \item \textbf{The team is responsible for the accuracy, completeness,
consistency, and appropriateness of the submitted document, regardless of
whether material was generated by a human or an AI system.}
\end{itemize}}

\newpage{}

\section*{Appendix --- Reflection}

\wss{Not required for CAS 741}

\input{../Reflection.text}

\begin{enumerate}
    \item Why is it important to create a development plan prior to starting the
    project?
    \item In your opinion, what are the advantages and disadvantages of using
    CI/CD?
    \item What disagreements did your group have in this deliverable, if any,
    and how did you resolve them?
\end{enumerate}

\newpage{}

\section*{Appendix --- Team Charter}

\wss{borrows from
\href{https://engineering.up.edu/industry_partnerships/files/team-charter.pdf}
{University of Portland Team Charter}}

\subsection*{External Goals}

\wss{What are your team's external goals for this project? These are not the
goals related to the functionality or quality fo the project.  These are the
goals on what the team wishes to achieve with the project.  Potential goals are
to win a prize at the Capstone EXPO, or to have something to talk about in
interviews, or to get an A+, etc.}

\subsection*{Attendance}

\subsubsection*{Expectations}

\wss{What are your team's expectations regarding meeting attendance (being on
time, leaving early, missing meetings, etc.)?}

\subsubsection*{Acceptable Excuse}

\wss{What constitutes an acceptable excuse for missing a meeting or a deadline?
What types of excuses will not be considered acceptable?}

\subsubsection*{In Case of Emergency}

\wss{What process will team members follow if they have an emergency and cannot
attend a team meeting or complete their individual work promised for a team
deliverable?}

\subsection*{Accountability and Teamwork}

\subsubsection*{Quality} 

\wss{What are your team's expectations regarding the quality
of team members' preparation for team meetings and the quality of the
deliverables that members bring to the team?}

\subsubsection*{Attitude}

\wss{What are your team's expectations regarding team members' ideas,
interactions with the team, cooperation, attitudes, and anything else regarding
team member contributions?  Do you want to introduce a code of conduct?  Do you
want a conflict resolution plan?  Can adopt existing codes of conduct.}

\subsubsection*{Stay on Track}

\wss{What methods will be used to keep the team on track? How will your team
ensure that members contribute as expected to the team and that the team
performs as expected? How will your team reward members who do well and manage
members whose performance is below expectations?  What are the consequences for
someone not contributing their fair share?}

\wss{You may wish to use the project management metrics collected for the TA and
instructor for this.}

\wss{You can set target metrics for attendance, commits, etc.  What are the
consequences if someone doesn't hit their targets?  Do they need to bring the
coffee to the next team meeting?  Does the team need to make an appointment with
their TA, or the instructor?  Are there incentives for reaching targets early?}

\subsubsection*{Team Building}

\wss{How will you build team cohesion (fun time, group rituals, etc.)? }

\subsubsection*{Decision Making} 

\wss{How will you make decisions in your group? Consensus?  Vote? How will you
handle disagreements? }

\end{document}